\documentclass{article}
\usepackage{graphicx} % Required for inserting images
\usepackage{booktabs}
\usepackage{amsmath}
\usepackage{amssymb}
\title{Results}
\author{srayan }
\date{March 2026}

\begin{document}

\maketitle



\section*{Single Solver}


\begin{table}[]
    \centering
    \resizebox{\textwidth}{!}{
    \begin{tabular}{|l|lrlrlrlr|}
    \toprule
    Methods & FinalError Poisson & P value Poisson & AUC Poisson & P value AUC Poisson & FinalError ConvDiff & P value ConvDiff & AUC ConvDiff & P value AUC ConvDiff \\
    \midrule 
    \multicolumn{9}{|c|}{Jacobi-related} \\
    \midrule
    Jacobi only & 0.496 (1.213) & 0.001152 & 1.064 (2.553) & 0.001070 & 0.176 (0.429) & 0.000971 & 0.407 (0.946) & 0.000902 \\
    HINTS-Jacobi & 0.761 (0.019) & 0.000000 & 0.463 (0.716) & 0.000000 & 0.454 (0.025) & 0.000000 & 0.235 (0.320) & 0.000011 \\
    Learned Greedy-Jacobi & 0.077 (0.216) & NaN & 0.223 (0.554) & NaN & 0.049 (0.149) & NaN & 0.129 (0.339) & NaN \\
    True-Greedy-Jacobi & 0.022 (0.018) & NaN & 0.113 (0.207) & NaN & 0.013 (0.014) & NaN & 0.060 (0.113) & NaN \\
    \midrule 
    \multicolumn{9}{|c|}{GS-related} \\
    \midrule
    GS only & 0.025 (0.060) & 0.001465 & 0.563 (1.350) & 0.001246 & 0.000 (0.000) & 0.038415 & 0.115 (0.266) & 0.000971 \\
    HINTS-GS & 0.724 (0.000) & 0.000000 & 0.382 (0.582) & 0.000000 & 0.456 (0.000) & 0.000000 & 0.150 (0.189) & 0.000000 \\
    Learned Greedy-GS & 0.003 (0.010) & NaN & 0.115 (0.285) & NaN & 0.000 (0.000) & NaN & 0.035 (0.095) & NaN \\
    True-Greedy-GS & 0.001 (0.001) & NaN & 0.067 (0.123) & NaN & 0.000 (0.000) & NaN & 0.022 (0.044) & NaN \\
    \midrule 
    \multicolumn{9}{|c|}{SymGS-related} \\
    \midrule
    SymGS only & 0.000 (0.000) & 0.016096 & 0.299 (0.725) & 0.001419 & 0.000 (0.000) & \textbf{0.789076} & 0.088 (0.204) & 0.001359 \\
    HINTS-SymGS & 0.709 (0.000) & 0.000000 & 0.296 (0.458) & 0.000000 & 0.460 (0.000) & 0.000000 & 0.128 (0.157) & 0.000000 \\
    Learned Greedy-SymGS & 0.000 (0.000) & NaN & 0.074 (0.194) & NaN & 0.000 (0.000) & NaN & 0.026 (0.050) & NaN \\
    True-Greedy-SymGS & 0.000 (0.000) & NaN & 0.041 (0.077) & NaN & 0.000 (0.000) & NaN & 0.019 (0.038) & NaN \\
    \bottomrule
    \end{tabular}}

    \caption{Single Solver result (Sym GS is Symmetric Gauss seidel)}
    \label{tab:placeholder}
\end{table}

\begin{figure}
    \centering
    \includegraphics[width=0.75\linewidth]{images2/deeponet_usage.png}
    \caption{DeepONet Usage}
    \label{fig:onetusage}
\end{figure}


\begin{figure}
    \centering
    \includegraphics[width=0.75\linewidth]{images2/solverusage_convdiff.png}
    \caption{Solver Usage ConvDiff}
    \label{fig:solverusage_conv}
\end{figure}

\begin{figure}
    \centering
    \includegraphics[width=0.75\linewidth]{images2/solverusage_poisson.png}
    \caption{Solver Usage Poisson}
    \label{fig:solverusage_pois}
\end{figure}

\begin{figure}
    \centering
    \includegraphics[width=0.75\linewidth]{images2/sample_error_convergence.png}
    \caption{Sample Error Convergence}
    \label{fig:error_convergence}
\end{figure}



\section{Multiple Solvers}

\begin{table}[]
    \centering
    \resizebox{\textwidth}{!}{
    \begin{tabular}{lllll}
    \toprule
    Methods & Final Error Poisson & AUC Poisson & Final Error ConvDiff & AUC ConvDiff \\
    \midrule
    Learned Greedy-SOR w/ 3 solvers & 0.000 (0.000) & 0.089 (0.214) & 0.000 (0.000) & 0.025 (0.057) \\
    True-Greedy-SOR w/ 3 solvers & 0.000 (0.000) & 0.048 (0.093) & 0.000 (0.000) & 0.013 (0.030) \\
    Learned Greedy-SOR w/ 4 solvers & 0.000 (0.000) & 0.061 (0.142) & 0.000 (0.000) & 0.011 (0.025) \\
    True-Greedy-SOR w/ 4 solvers & 0.000 (0.000) & 0.037 (0.076) & 0.000 (0.000) & 0.009 (0.022) \\
    Learned Greedy-SOR w/ 5 solvers & 0.000 (0.000) & 0.049 (0.115) & 0.000 (0.000) & 0.013 (0.031) \\
    True-Greedy-SOR w/ 5 solvers & 0.000 (0.000) & 0.035 (0.079) & 0.000 (0.000) & 0.009 (0.022) \\
    Learned Greedy-SOR w/ 6 solvers & 0.000 (0.000) & 0.053 (0.123) & inf (nan) & inf (nan) \\
    True-Greedy-SOR w/ 6 solvers & 0.000 (0.000) & 0.035 (0.079) & 0.000 (0.000) & 0.009 (0.022) \\
    \bottomrule
    \end{tabular}}
    \caption{Caption}
    \label{tab:placeholder}
\end{table}





\begin{figure}
    \centering
    \includegraphics[width=0.75\linewidth]{images2/multiple_solverusage_poisson.png}
    \caption{Multiple Solver Usage Poisson}
    \label{fig:multiplesolver}
\end{figure}

\begin{figure}
    \centering
    \includegraphics[width=0.75\linewidth]{images2/multiple_solverusage_convdiff.png}
    \caption{Multiple solver usage Convdiff}
    \label{fig:placeholder}
\end{figure}

\section{Notes}

Single Solver suggestion
\begin{itemize}
    \item why are the row labels bold?
    \item Title for solver usage needs the equation in it
    \item poisson jacobi and conv diff jacobi error convergence plots look suspiciously similar. lol they aren't
    \item I need to retrain greedy SOR with 6 solvers for ConvDiff
    \item the multiple solver usage plots need better titles
    \item Table 1
    \begin{itemize}
        \item defer p values and true greedy stuff to the appendix 
        \item bold the best one
        \item include explanation of the bolded p value in the appendix
        \item include two more solvers - appenidx work, SOR and weighted jacobi
    \end{itemize}
    \item Figure 1
    \begin{itemize}
        \item include the true greedy and learned greedy comparison in the appendix
        \item we want the learned greedy deeponet usage in the main body
        \item we want to to point out two things:
        \begin{itemize}
            \item it is using the deeponet
            \item using it diffently across samples
        \end{itemize}
        \item cut off the iteration sooner so maybe at 100 or 150
        \item move the true greedy comparisons to the appendix
    \end{itemize}
    \item Figure 2
    \begin{itemize}
        \item move this to the appendix
    \end{itemize}
    \item New Figure: final prediction from HINTS, solver only, learned greedy and ground truth for each solver (jacbi, gs, symgs) - 12 heatmaps
    \item Table 2
    \begin{itemize}
        \item include the final error and Auc for sor only, sor 1.25, etc lol
        \item we want two convey things:
        \begin{itemize}
            \item monotonc decreasing trend as we increase the numvr solvers
            \item and it's just better than having each of them run indivdually for 300 iteration
        \end{itemize}
        \item move the true greedy comparisons to the appendix
    \end{itemize}
    \item Figure 5:
    \begin{itemize}
        \item router is using the entire suite of solvers
        \item moving to the appendix
        \item highlight that this is beneficial to learn a sample-dependent router
    \end{itemize}
    \item mode wise errors also we need
    
\end{itemize}


In this section, we empirically demonstrate the fast and uniform convergence of the approximate greedy router on Poisson and convection–diffusion (ConvDiff) equations posed on the unit domain $D = [0,1]^2$.

For Poisson, we consider the constant-coefficient model
\begin{equation*}
    - \Delta u(\mathbf{x}) = f(\mathbf{x}), \quad \mathbf{x}\in D,
\end{equation*}
and for ConvDiff, we adopt the sign convention
\begin{equation*}
    -\Delta u(\mathbf{x}) + 20\,\partial_{x_1}u(\mathbf{x}) + 20\,\partial_{x_2}u(\mathbf{x}) = f(\mathbf{x}), \quad \mathbf{x}\in D.
\end{equation*}
We impose periodic boundary conditions for both equations, and center $f$ to satisfy the compatibility condition $\int_D f(\mathbf{x})\,d\mathbf{x} = 0$. The domain is discretized on a $31\times31$ uniform grid.

Data is sampled from a zero-mean hierarchical Gaussian random field on the periodic domain with covariance operator $\alpha(-\Delta + \beta I)^{-\gamma}$, where $\alpha,\beta,\gamma$ vary across samples (see \Cref{sec:training details}).

We report performance over $64$ test samples using the mean final error $\|e_h^{(T)}\|_2$ and the mean area under the curve (AUC), defined as $\sum_{t=1}^T \|e_h^{(t)}\|_2$. While the final error measures terminal accuracy, AUC captures convergence over the full trajectory and penalizes non-monotone behavior.

\paragraph{Comparing Greedy with HINTS.}
We consider Jacobi, Gauss–Seidel (GS), and symmetric Gauss–Seidel (SymGS) solvers, together with a DeepONet model. Baselines include single-solver schedules and HINTS variants (HINTS-Jacobi, HINTS-GS, HINTS-SymGS), where the DeepONet correction is applied every 24 iterations. 

We train LSTM-based routers using the surrogate loss in \Cref{eq:routersurrogate} for three solver-access sets: Jacobi+DeepONet (Greedy-Jacobi), GS+DeepONet (Greedy-GS), and SymGS+DeepONet (Greedy-SymGS). We use an LSTM to capture the sequential nature of routing decisions, where the benefit of each action depends on the evolving error trajectory.

Comparisons are restricted to methods with identical solver access. Accordingly, results are grouped by solver family (Jacobi, GS, and SymGS). Each method is run for 300 iterations.


As shown in \Cref{tab:errorcomparison}, the greedy router achieves the lowest final error and mean AUC across all solver families: Greedy-Jacobi, Greedy-GS, and Greedy-SymGS consistently outperform both single-solver and HINTS baselines. While HINTS improves over fixed schedules, it exhibits sawtooth error trajectories (\Cref{fig:errorcomparison}) due to periodic application of the DeepONet even when suboptimal. In contrast, the greedy routers select solvers that yield immediate error reduction, resulting in lower AUC and near-monotone convergence.

The learned routing policy is highly non-uniform across samples (\Cref{fig:deeponet}): the DeepONet correction is often unused after approximately 40 iterations for many instances. This highlights the need for sample-dependent routing. HINTS, by design, continues to invoke the neural operator at fixed intervals regardless of its effectiveness, leading to higher error and degraded AUC.

\end{document}